% QMB_n08_qubit_formalismus.tex — Kapitel 8: Qubit-Formalismus
% Inhaltliche Grundlage: Dok. 175 (März 2026), Dok. 034, Dok. 230; eigene Darstellung
\chapter{Mehr als zwei Zustände: das Qubit geometrisch}
\label{ch:qubit}

\section{Was steckt zwischen $|0\rangle$ und $|1\rangle$?}

Spin-up und Spin-down sind zwei Punkte. Zwei Punkte sind eindeutig
beschriftet --- ein Bit Information. Aber ein Qubit ist kein Bit. Die
eigentliche Kraft des Qubits liegt nicht darin, dass es $|0\rangle$
\emph{oder} $|1\rangle$ sein kann; das kann eine klassische Münze auch.
Sie liegt darin, dass es beides in einem kontinuierlichen Verhältnis
überlagern kann:
\begin{equation}
  |\psi\rangle = \cos\tfrac{\theta}{2}\,|0\rangle
  + e^{i\phi}\sin\tfrac{\theta}{2}\,|1\rangle.
\end{equation}
Zwei Parameter, der Mischungswinkel $\theta\in[0,\pi]$ und die relative
Phase $\phi\in[0,2\pi)$, parametrieren eine Kugeloberfläche: die
Bloch-Kugel. Auf ihr gibt es unendlich viele Punkte --- nicht nur die
Pole $|0\rangle$ und $|1\rangle$, sondern den gesamten Äquator, alle
Breitenkreise, jeden Punkt dazwischen. Jeder ist ein gültiger,
physikalisch unterscheidbarer Zustand.

\begin{keypoint}[Ein Qubit kodiert nicht 2 Zustände, sondern unendlich viele]
Die zwei klassischen Pole sind Sonderfälle einer kontinuierlichen
Familie. Was ein Qubit vom Bit unterscheidet, ist nicht die Zahl der
diskreten Zustände, sondern die kontinuierliche Geometrie des
Zustandsraums. Das Auslesen koppelt an die Konfiguration und treibt sie
in den nächstgelegenen stabilen Extremzustand. Das Ergebnis ist in
FFGFT-Sicht deterministisch --- es hängt vom genauen Ausgangszustand
ab, den wir vor der Messung nicht vollständig kennen. Die Zustände
selbst sind diskret; die Kontinuität liegt im Phasenraum \emph{vor}
der Messung.
\end{keypoint}

\section{Der zylindrische Phasenraum}

Die Bloch-Kugel ist die Standarddarstellung. Die FFGFT hat eine
physikalisch direktere Alternative: den zylindrischen Phasenraum mit
drei Koordinaten $(z,r,\theta)$. Dabei ist $z\in[-1,+1]$ die Höhe
entlang der Zylinderachse --- der Anteil der Torsionskonfiguration in
Richtung $|0\rangle$ gegenüber $|1\rangle$; $r\in[0,1]$ der Radius im
Querschnitt --- die Stärke der intermediären Torsion, also wie weit der
Zustand von beiden Polen entfernt ist; $\theta\in[0,2\pi)$ der
Azimutwinkel --- die relative Phase zwischen den Basiskomponenten. Die
Beziehung zur Bloch-Kugel ist direkt: $z = \cos\theta_{\mathrm{Bloch}}$,
$r = |\sin\theta_{\mathrm{Bloch}}|$; die formale Brücke steht in
Kapitel~\ref{ch:hilbert}.

Der Zylinder ist kein geometrisches Schönheitsmittel: Er macht
Gatteroperationen zu geometrischen Transformationen im Realraum.

\begin{center}
\small\renewcommand{\arraystretch}{1.4}
\begin{tabular}{lp{2.4cm}p{4.2cm}}
\toprule
\textbf{Gatter} & \textbf{Matrix} & \textbf{Zylindertransformation} \\
\midrule
Hadamard $H$ & $\frac{1}{\sqrt2}\begin{pmatrix}1&1\\1&-1\end{pmatrix}$ & $z\leftrightarrow r$ vertauschen, Phase um $\pi/2$ drehen: bewegt Pol auf Äquator \\
Phase $Z$ & $\begin{pmatrix}1&0\\0&-1\end{pmatrix}$ & $\pi$ zur Phase $\theta$ addieren: Rotation um die Zylinderachse \\
Bit-Flip $X$ & $\begin{pmatrix}0&1\\1&0\end{pmatrix}$ & $(z,r)$-Drehung um $\alpha = \pi K_{\mathrm{frak}}$ \\
\bottomrule
\end{tabular}
\end{center}

Im zylindrischen Formalismus ist jede Gatteroperation eine geometrische
Transformation des Torsionsvektors. Das Hadamard-Gatter dreht keine
Matrix --- es kippt die Konfiguration von der Längsachse auf die
Transversalebene. Das $Z$-Gatter rotiert die Windung um ihre eigene
Achse. Der Faktor $K_{\mathrm{frak}}$ im $X$-Gatter ist kein nachträglich
eingefügter Korrekturfaktor; er folgt aus der fraktalen
Raumzeitgeometrie und bedeutet: Jede Qubit-Rotation ist leicht kleiner
als $\pi$, weil die Raumzeit selbst minimal von der flachen Geometrie
abweicht. Das ist eine prüfbare Vorhersage: Alle $\pi$-Gatter sollten
systematisch um $\Delta\alpha\approx0{,}013\,\%$ vom idealen Wert
abweichen.\status{S}

\section{Quantenregister: exponentielles Wachstum durch Kombination}

Ein einzelnes Qubit hat unendlich viele Zwischenzustände, liefert beim
Auslesen aber immer einen der zwei Pole. Die eigentliche Stärke kommt
durch Kombination.

\textbf{Produktzustände: additive Klassik.} Zwei unabhängige Qubits
haben zusammen vier Basiszustände. Der allgemeine Zustand ohne
Verschränkung ist $(\alpha_A|0\rangle+\beta_A|1\rangle)\otimes
(\alpha_B|0\rangle+\beta_B|1\rangle)$ --- zwei getrennte Bloch-Kugeln,
kein Gewinn über zwei unabhängige Qubits. Der Zustandsraum ist additiv:
$2+2 = 4$ Parameter.

\textbf{Verschränkte Zustände: multiplikatives Wachstum.} Sobald die
Qubits verschränkt sind, kann der gemeinsame Zustand nicht mehr als
Produkt geschrieben werden. Der allgemeine Zweiqubitzustand
\begin{equation}
  |\Psi\rangle = \alpha|00\rangle+\beta|01\rangle+\gamma|10\rangle+\delta|11\rangle
\end{equation}
braucht vier komplexe Amplituden, also $2\cdot4-2 = 6$ reelle Parameter.
Verschränkte Qubits tragen gemeinsame Information, die in keinem der
Einzelqubits steckt. Für $n$ Qubits wächst der Hilbertraum exponentiell,
$d(n) = 2^n$, mit $2\cdot2^n-2$ reellen Parametern:

\begin{center}
\small
\begin{tabular}{rrrl}
\toprule
$n$ Qubits & Basiszustände & Reelle Parameter & Klassisch \\
\midrule
1 & 2 & 2 & 1 Bit \\
2 & 4 & 6 & 2 Bits \\
3 & 8 & 14 & 3 Bits \\
10 & 1024 & 2046 & 10 Bits \\
50 & $10^{15}$ & $2\times10^{15}$ & 50 Bits \\
300 & $10^{90}$ & $2\times10^{90}$ & $\approx$ Atome im Universum \\
\bottomrule
\end{tabular}
\end{center}

Das ist der eigentliche Quantenvorteil: Nicht ein einzelnes Qubit ist
mächtig, sondern das kollektive Wachstum eines verschränkten Registers.

\textbf{Bell-Zustände.} Die vier maximal verschränkten
Zweiqubitzustände bilden eine Orthonormalbasis des verschränkten
Unterraums:
\begin{align}
  |\Phi^\pm\rangle &= \tfrac{1}{\sqrt2}(|00\rangle\pm|11\rangle),\\
  |\Psi^\pm\rangle &= \tfrac{1}{\sqrt2}(|01\rangle\pm|10\rangle).
\end{align}
Misst man Qubit~A, ist das Ergebnis von Qubit~B festgelegt, egal wie
weit entfernt. Im FFGFT-Bild ist ein Bell-Zustand keine Fernwirkung
zwischen zwei Qubits, sondern eine topologisch verbundene
Torsionskonfiguration über zwei Resonatoren: Beide teilen einen
gemeinsamen Torsionsknoten, dessen Auflösung bei jeder lokalen Messung
beide Enden zugleich bestimmt. Das ist kein Signal mit
Überlichtgeschwindigkeit, sondern ein gemeinsamer geometrischer Zustand,
der bei der Präparation etabliert wurde (Kapitel~\ref{ch:verschraenkung}).
Die frühere Dämpfungsformel für die Bell-Korrelation, die den CHSH-Wert
um $\sim10^{-4}$ absenkte, ist durch die Zwei-Observablen-Lesart aus
Kapitel~\ref{ch:chsh_exp} ersetzt.

\textbf{Hardware-Validierung.} Die FFGFT-Vorhersagen für Bell-Zustände
wurden auf IBM Brisbane und IBM Sherbrooke (je 127 Qubits, Heavy-Hex-
Gitter) geprüft: Bell-Zustands-Fidelität $97{,}17\,\%$; FFGFT-Algorithmus
und Standard-QM stimmen innerhalb der Messgenauigkeit überein;
Wiederholvarianz $10^{-4}$ bis $10^{-3}$; Abweichungen in $P(00)$ und
$P(11)$ von $1$--$5\,\%$ sind Hardware-Rauschen.\status{K} Die
$\xi$-Abweichung von der idealen Bell-Korrelation liegt unterhalb der
Rauschschwelle und konnte nicht isoliert werden. Der vollständige
Python-Code ist im Repositorium verfügbar.

\section{Drei Wege zu mehr Zuständen}

Aus der Resonatorphysik ergeben sich drei grundlegend verschiedene
Strategien, mehr als zwei Zustände zu kodieren:

\begin{center}
\small\renewcommand{\arraystretch}{1.3}
\begin{tabular}{p{1.6cm}p{2.3cm}p{1.6cm}p{2.5cm}}
\toprule
\textbf{Strategie} & \textbf{Mehr Zustände durch} & \textbf{Wachstum} & \textbf{FFGFT-Bedeutung} \\
\midrule
Superposition (1 Qubit) & kontinuierliche Überlagerung & $\infty$, aber 1 Bit bei Messung & intermediäre Torsion \\
Register ($n$ Qubits) & Verschränkung & $2^n$ & topologisch verbundene Knoten \\
Qudit ($d>2$) & höhere Niveaus, Orbitale & $d$ pro Objekt & $d$ stabile Konfigurationen in einem Resonator \\
\bottomrule
\end{tabular}
\end{center}

Keine Strategie dominiert überall. Registerverschränkung skaliert am
stärksten, braucht aber $n$ physische Objekte und fehlerkorrigierte
Verschränkungsoperationen. Qudit-Kodierung gewinnt Information pro Objekt
($\log_2d>1$); zwei Ququarts ($d = 4$) tragen $4^2 = 16$ Zustände, vier
Qubits $2^4 = 16$ --- mathematisch äquivalent, physikalisch verschieden
robust. Superposition ist immer da: Sie ist die Grundlage der
Interferenz, der eigentlichen Triebkraft von Quantenalgorithmen.

\section{Quanteninterferenz: warum Superposition nützlich ist}

Superposition allein wäre eine bloße Überlagerung; das Auslesen lieferte
immer einen der zwei Pole. Was Superposition zur Ressource macht, ist die
Interferenz: Amplituden können sich konstruktiv verstärken oder
destruktiv auslöschen, je nach Phase.\status{E} Das ist das Prinzip
hinter allen nützlichen Quantenalgorithmen: Man präpariert das Register
in einer gleichmäßigen Zwischentorsion über alle Eingaben (Hadamard);
wendet eine Funktion an, die die Phasen der richtigen Antworten erhöht
und die der falschen erniedrigt; lässt Interferenz die falschen
Antworten in instabile Konfigurationen treiben und die richtige in eine
stabile; und das Auslesen treibt das System in den nächsten stabilen
Pol, der nach der Interferenz mit überwältigender Häufigkeit der
gesuchten Antwort entspricht. Ohne Interferenz: Monte Carlo. Mit
Interferenz: Grover ($\sqrt N$ statt $N$ Schritte), Shor (Faktorisierung
in polynomialer Zeit).

In der FFGFT ist die Phase keine abstrakte komplexe Zahl, sondern die
Rotationsstellung der intermediären Torsionskonfiguration. Konstruktive
Interferenz: Zwei Torsionsphasen liegen gleich und verstärken sich zu
einer geometrisch stabileren Konfiguration. Destruktive Interferenz:
Zwei entgegengesetzte Phasen löschen sich aus; die resultierende
Konfiguration hat keinen stabilen geometrischen Anker und wird beim
Auslesen nicht angesteuert. Ein Quantenalgorithmus ist in dieser
Sprache eine Folge von Torsionsrotationen, die die Phasenkonfiguration
des Registers so ausrichtet, dass die gesuchte Antwort in einer stabilen
Konfiguration endet.

\begin{laierspur}
Wasserwellen, die sich treffen, verstärken sich, wenn Berg auf Berg
trifft, und löschen sich aus, wenn Berg auf Tal trifft. Ein
Quantenrechner steuert Wahrscheinlichkeiten wie Wellen: Er lässt die
richtigen Antworten wachsen und die falschen verschwinden. Das geht nur,
weil Quantenzustände eine Phase haben, die klassischen Bits fehlt.
\end{laierspur}

\section{QND-Messung: Auslesen ohne Zerstören}

Die Auslesewechselwirkung treibt die intermediäre Konfiguration in
einen Pol --- für invasive Messverfahren. Es gibt eine Klasse von
Verfahren, die diesen Schritt nicht vollziehen: die
Quanten-Nichtdemolitionsmessung (QND). Sie liest eine Observable $\hat A$
aus, ohne den Eigenzustand zu verändern; die formale Bedingung ist
$[\hat A,\hat H_{\mathrm{Messung}}] = 0$.\status{E} Praktisch: Derselbe
Zustand kann wiederholt gemessen werden und liefert jedes Mal dasselbe
Ergebnis.

Das wichtigste Beispiel ist die Faraday- oder Kerr-Rotation: Ein linear
polarisierter Laserstrahl passiert den Quantenpunkt ohne Absorption und
dreht seine Polarisationsebene um $\theta_F\propto S_z$. Kein Photon wird
absorbiert, kein Energieübertrag in die Spinmode, der Spinzustand ist
vor und nach der Messung identisch.

\textbf{QND als empirischer Beleg für Determinismus.} Das Argument ist
direkt: Die Faraday-Rotation misst $S_z$ ohne den Zustand zu verändern;
wiederholte Messungen liefern immer dasselbe Ergebnis; also hatte der
Zustand diesen Wert vor jeder Messung, definit, nicht als abstrakte
Superposition, die erst durch Messung Realität wird. In der FFGFT ist
die Faraday-Rotation zustandserhaltend, weil eine stabile
Torsionskonfiguration durch eine dispersive, nichtresonante Kopplung
nicht aus ihrem Zustand getrieben werden kann. Das Photon sieht die
Torsion --- es dreht seine Polarisation entsprechend ---, überträgt aber
keine Energie, weil es nicht resonant ist.\status{S}

\begin{center}
\small\renewcommand{\arraystretch}{1.3}
\begin{tabular}{p{2.4cm}p{3cm}p{2.8cm}}
\toprule
\textbf{Verfahren} & \textbf{Wechselwirkung} & \textbf{Zustand danach} \\
\midrule
Resonante Fluoreszenz & Photon absorbiert, regt Übergang an & verändert \\
Spin-zu-Ladungs-Konversion & Elektron tunnelt aus dem Resonator & zerstört \\
Faraday/Kerr & Photon passiert dispersiv & unverändert: QND \\
\bottomrule
\end{tabular}
\end{center}

Die Standardlesart, der Zustand existiere erst durch Messung, ist mit
QND-Messungen schwer vereinbar. In der FFGFT ist die
Torsionskonfiguration eine geometrische Realität: stabil, definit,
präexistent. Die Messung liest sie ab; sie erschafft sie nicht.

\textbf{Grenzen der QND.} QND ist möglich für Observablen, die mit dem
Messoperator kommutieren, typisch $S_z$. Eine intermediäre Konfiguration
--- ein Zustand, der kein $S_z$-Eigenzustand ist --- wird durch eine
$S_z$-QND-Messung in einen Eigenzustand getrieben, deterministisch
abhängig vom Ausgangszustand; die Zwischentorsion ist danach weg.

\section{Die Potentiallandschaft}

Das vollständige Bild der Qubit-Zustände ist eine Potentiallandschaft
mit zwei stabilen Minima und einem Sattel dazwischen. Das Minimum
$|\uparrow\rangle$ ist eine stabile Konfiguration, tief genug, dass
schwache Kopplungen den Zustand nicht herausheben --- QND möglich. Das
Minimum $|\downarrow\rangle$ ist spiegelbildlich gleich tief. Der Abhang
ist die intermediäre Torsion: Die Phase bestimmt, wie weit der Zustand
von welchem Minimum entfernt ist; kein echtes Minimum, metastabil auf
der Kohärenzzeit $T_2$. Jeder Punkt dazwischen ist ein geometrisch
definierter Zustand, kontinuierlich parametriert, diskret in den
Endpunkten.

Die Auslesestärke entscheidet: Ein Zustand im stabilen Minimum bleibt
bei schwacher Kopplung (QND) und kann bei starker Kopplung herausgehoben
werden (invasiv). Ein Zustand auf dem Abhang rollt schon bei schwacher
Kopplung ins nächste Minimum --- welches, bestimmt die genaue Phase,
epistemisch unbekannt, deterministisch; bei starker Kopplung dasselbe,
nur schneller.

Die entscheidende Einsicht: Es gibt keine prinzipielle Grenze zwischen
\enquote{QND-möglich} und \enquote{QND-unmöglich}, nur stabilere und
weniger stabile Zustände in einer kontinuierlichen Landschaft. Der
Sprung in die Endposition ist nicht zufällig: Ein Zustand nahe
$|\uparrow\rangle$ landet in $|\uparrow\rangle$, einer nahe
$|\downarrow\rangle$ in $|\downarrow\rangle$, einer auf dem Sattel
abhängig von der feinen Phase. Die scheinbare Zufälligkeit der
Standard-QM ist die epistemische Unkenntnis dieser Phase. Das Bild gilt
für alle $d$ Niveaus eines Qudits: $d$ Minima mit $d-1$ Sattelpunkten.

\section{Zwei Ebenen der $\xi$-Hierarchie: Spin-Qubit und photonischer Resonator}

Die bisherigen Abschnitte haben das Qubit am Spin-Qubit entwickelt ---
einem Elektron in einem Nanometerkäfig. Es gibt eine qualitativ andere
Realisierungsklasse: den photonischen Resonator. Beide erzwingen
diskrete Zustände durch Geometrie, sitzen aber auf verschiedenen Stufen
der $\xi$-Hierarchie, mit Konsequenzen für Temperatur, Stabilität und
die Art der Information.

\textbf{Spin-Qubit: minimaler Resonator auf Stufe $N\approx8$.} Ein
einzelnes Elektron, ein Energieniveau, zwei Orientierungen. Der Abstand
der Spinzustände in einem typischen Feld von 1\,T ist
$\Delta E_Z = g^*\mu_BB\approx0{,}1$--$1$\,meV; die thermische Energie
bei Raumtemperatur ist $k_BT = 26$\,meV. Also $\Delta E_Z\ll k_BT$:
Thermisches Rauschen wirft das System ständig aus dem Minimum. Tiefe
Temperaturen unter 1\,K sind nicht optional, sondern zwingend --- nicht
um Superposition zu erhalten, sondern um überhaupt in einem der Minima
zu bleiben.

\textbf{Photonischer Resonator: kollektive Geometrie auf Stufe
$N\approx9$--$10$.} Ein Wellenleiter aus Lithiumniobat, ein
Mach-Zehnder-Interferometer, ein Ringresonator --- Resonatoren aus
$10^{18}$ bis $10^{23}$ Atomen in Gitterordnung. Der Resonator ist nicht
ein einzelnes Teilchen, sondern die makroskopische Geometrie des ganzen
Kristalls; das Photon gehört der kollektiven Mode. Die Diskretheit kommt
aus den Randbedingungen der makroskopischen Geometrie. Die Energie eines
optischen Photons ist $1$--$2$\,eV, also $E/k_BT\approx40$: Thermisches
Rauschen kann keine optischen Photonen erzeugen. Photonik funktioniert
bei Raumtemperatur nicht, weil sie robuster gebaut ist, sondern weil
sie auf einer höheren $\xi$-Stufe operiert: Die geometrische Robustheit
ist durch den Einkristall mit $10^{23}$ Gitterplätzen gegeben, die
thermische durch die eV-Skala --- vier Größenordnungen Abstand.

\textbf{Nicht auf zwei Zustände beschränkt.} In einem Mach-Zehnder-
Interferometer läuft das Photon durch zwei Pfade mit kontinuierlich
einstellbarer Phase; das Interferenzergebnis $\cos^2(\phi/2)$ ist eine
kontinuierliche Übertragungsfunktion. Ein photonischer Resonator kann
beide Modi betreiben: digital mit $\phi\in\{0,\pi\}$ und Transmission
1 oder 0, oder analog mit kontinuierlicher Phase --- unendlich viele
Zwischenzustände, alle bei Raumtemperatur stabil. Beim Spin-Qubit liegt
die analoge Information in einem thermisch instabilen Zwischenzustand,
$T_2$ bricht bei Raumtemperatur in Nanosekunden zusammen; beim
photonischen Resonator liegt sie in der Phase des Photons, und die ist
thermisch stabil.

\begin{center}
\small\renewcommand{\arraystretch}{1.3}
\begin{tabular}{p{2.6cm}p{2.8cm}p{3cm}}
\toprule
 & \textbf{Spin-Qubit} & \textbf{Photonischer Resonator} \\
\midrule
$\xi$-Stufe & $N\approx8$ (nm) & $N\approx9$--10 ($\mu$m--mm) \\
Resonator & 1 Elektron im Käfig & $10^{18}$--$10^{23}$ Atome kollektiv \\
Energieskala & $0{,}1$--1\,meV & 1--2\,eV \\
$\Delta E/k_BT$(300\,K) & $0{,}004$--$0{,}04$ & $\approx40$ \\
Kühlung nötig? & ja, $T<1$\,K & nein \\
Analoge Zustände & fragil ($T_2\sim$ns) & stabil \\
Diskretheit aus & Confinement (1 Teilchen) & Modengeometrie \\
Typische Anwendung & digitales Qubit & analoge Signalverarbeitung, photonisches Rechnen \\
\bottomrule
\end{tabular}
\end{center}

Spin-Qubit und photonischer Resonator sind in der FFGFT keine
grundlegend verschiedenen Systeme, sondern Resonatoren auf verschiedenen
Stufen derselben Skalierungsleiter. Jede Stufe trägt die Energie um den
Faktor $1/\xi\approx7500$ höher als die vorherige; von meV (Stufe 8) zu
eV (Stufe 9--10) ist der Faktor $10^3$--$10^4$, konsistent mit der
Skalierung.\status{S} Das erklärt auch, warum biologische Systeme
photonische Mechanismen (Photosynthese, Vogelnavigation) und keine
Spin-Qubits für Quanteneffekte nutzen: Die Natur arbeitet auf der
thermisch robusten Stufe.

\section{Die analoge Auflösungsgrenze}

Der photonische Resonator ist thermisch robust, aber nur, solange man
nicht zu fein auflöst. Wer analoge Zwischenwerte unterscheiden will,
teilt den Energieraum in $N$ Stufen mit $\Delta E_{\mathrm{Stufe}} =
E_{\mathrm{Photon}}/N$, und die Bedingung $\Delta E_{\mathrm{Stufe}}\gg
k_BT$ ergibt $N\ll E_{\mathrm{Photon}}/k_BT$. Bei Raumtemperatur und
optischen Photonen: $N_{\max}\approx40$, entsprechend
$\log_2 40\approx5{,}3$\,Bit pro Photon. Das photonische System entkommt
dem thermischen Problem nicht --- es verschiebt es auf eine feinere
Skala. Mehr Auflösung erfordert Kühlung oder eine noch höhere
Energieskala.

Das ist die älteste Ingenieursweisheit der Nachrichtentechnik,
geometrisch begründet: Die Hochfrequenztechnik klettert seit hundert
Jahren dieselbe Hierarchie hoch. Bei jeder Frequenzstufe verbessert sich
das Signal-Rausch-Verhältnis, weil $E/k_BT$ mit jeder Stufe wächst.

\begin{center}
\footnotesize\renewcommand{\arraystretch}{1.25}
\begin{tabular}{p{2.2cm}p{1.3cm}p{1.6cm}p{1.2cm}p{2.9cm}}
\toprule
\textbf{Technologie} & \textbf{Frequenz} & \textbf{$E/k_BT$} & \textbf{$N$-Stufe} & \textbf{Regime} \\
\midrule
AM-Rundfunk & 1\,MHz & $1{,}5\times10^{-7}$ & $\approx5{,}5$ & thermisch dominiert \\
UKW & 100\,MHz & $1{,}5\times10^{-5}$ & $\approx6{,}4$ & rauschbegrenzt \\
Mikrowelle & 10\,GHz & $1{,}5\times10^{-3}$ & $\approx7{,}3$ & beherrschbar \\
5G/mmWave & 30\,GHz & $5\times10^{-3}$ & $\approx7{,}5$ & Übergang \\
THz & 1\,THz & $0{,}15$ & $\approx7{,}8$ & $E\sim k_BT$: schwierig \\
Nahinfrarot & 200\,THz & 31 & $\approx8{,}8$ & thermisch robust \\
Glasfaser & 300\,THz & 46 & $\approx9{,}0$ & $N_{\max}\approx46$ \\
UV & 1000\,THz & 154 & $\approx9{,}4$ & $N_{\max}\approx150$ \\
\bottomrule
\end{tabular}
\end{center}

Von $N = 6$ (AM) nach $N = 9$ (Optik) überbrücken drei Stufen den Faktor
$(1/\xi)^3\approx4\times10^{11}$; tatsächlich liegt die Glasfaserfrequenz
$10^8$-fach über AM-Rundfunk. Die Glasfaser ist in diesem Bild keine
neue Technologie, sondern die Fortsetzung desselben Trends auf die
nächste Stufe. Und die Auflösungsgrenze gilt auf jeder Stufe: Wer bei
optischen Wellenlängen mehr als 40 analoge Stufen auflösen will, muss
kühlen --- wie der Radioingenieur 1930 seinen Empfänger abschirmte. Die
universelle Grenze lautet
\begin{equation}
  N_{\max}(T,N_\xi)\approx\frac{E_0\,(1/\xi)^{N_\xi-N_0}}{k_BT}.
\end{equation}

\textbf{Oberhalb der Optik.} Die Optik ist nicht das Ende der Leiter,
aber das Ende der nutzbaren Zone für geführte photonische Übertragung.
Im nahen UV (3--10\,eV, $N\approx9{,}3$) beginnt Photoionisation: Der
Resonator wird durch sein eigenes Signal beschädigt. Bei weichem Röntgen
(100\,eV--1\,keV, $N\approx10$) nähert sich die Wellenlänge dem
Gitterabstand, Bragg-Streuung an einzelnen Gitterebenen, die
Wellenleitung bricht zusammen. Bei hartem Röntgen fliegt das Photon
durch die Gitterlücken, Materie wird teils transparent, es gibt keinen
Brechungsindex und keinen Resonator. Bei Gammaenergien wechselwirkt das
Photon nur noch mit Atomkernen --- eine andere $\xi$-Stufe. Oberhalb
1\,MeV ($N\approx11{,}3$) setzt Paarerzeugung ein: die natürliche
Obergrenze. In der FFGFT gibt es keine Kanonenkugeln und keine Wellen als
getrennte Konzepte, nur Torsionskonfigurationen auf verschiedenen Skalen;
Wellenleitung ist eine Skalenresonanz zweier Torsionsmuster, des
Kristalls auf Stufe 7--8 und des Photons auf Stufe 9, und sie bricht
zusammen, wenn die Skalen inkompatibel werden.

\begin{keypoint}[Was dieses Kapitel festhält]
Das Qubit ist eine Mode im zylindrischen Phasenraum; Gatter sind
Torsionsrotationen; Register wachsen exponentiell durch topologische
Verbindung; Interferenz ist Phasenüberlagerung von Torsionen; QND-Messung
belegt die Präexistenz stabiler Pole; die Potentiallandschaft macht den
Messausgang deterministisch bei epistemisch unbekannter Phase. Spin-Qubit
und Photonik sind Resonatoren auf verschiedenen $\xi$-Stufen, und die
thermische Auflösungsgrenze gilt auf jeder Stufe gleich.
\end{keypoint}

\quelle{Dok.~175 (Qubit-Zustandsräume, März 2026), Dok.~034 (zylindrischer Formalismus), Dok.~173 (Resonatorbedingungen), Dok.~174 (Torsion, Faraday-Auslese), Dok.~230}
